\chapter{Conclusiones y trabajo futuro}
\label{chap:conclusiones}

% Versión 4: incorpora las conclusiones del OE2 y del OE3 acotadas según la síntesis de
% la dirección (revisión EF-01-C2 y curvas de exposición, rama principal 08cdc0c), una
% conclusión general y el trabajo futuro del modelo.
% Versión 3: conclusión del OE1 según SINTESIS_RESULTADOS_OE1.md (revisión 4529c49).

Este capítulo resume lo que los resultados permiten afirmar respecto de cada objetivo específico y señala el trabajo que queda abierto. Contiene una conclusión por objetivo, una conclusión general que las articula con el objetivo general del proyecto y las líneas de trabajo futuro.

\section{Conclusión del primer objetivo específico}
\label{sec:conc-oe1}

Se obtuvo un conjunto de datos sintéticos de imágenes muográficas por transmisión, con diferentes niveles de estadística de conteo, generado mediante simulación Monte Carlo del transporte de muones primarios en Geant4 a través de diecinueve escenarios volcánicos idealizados. El conjunto contiene 114 referencias, seis por escenario, obtenidas cada una con $3{,}44\times10^{8}$ primarios, y 9\,120 pares de entrada y referencia producidos por adelgazamiento binomial con fracciones de 0,50, 0,10, 0,05 y 0,01, particionados por geometría en doce escenarios de entrenamiento, tres de validación, tres de prueba y uno de generalización. Las 114 referencias resultaron completas e íntegras, la respuesta en las 1\,710 regiones de borde quedó dentro de la tolerancia fijada de antemano, y 107 de las 114 vistas alcanzaron la meta descriptiva de suficiencia del conteo en roca, mientras que las siete restantes quedaron a una o dos regiones de ella y se conservaron conforme al acuerdo previo. El posproceso desde los histogramas y la construcción de los pares se reprodujeron de forma independiente con contenido idéntico, y la copia de distribución se recuperó y verificó completa. Con ello, los resultados sustentan la obtención del conjunto sintético que el OE1 formuló, bajo el alcance idealizado adoptado: conteos direccionales sintéticos de transmisión de una familia de edificios cónicos con conductos, cavidades y cuerpos densos.

Esa conclusión tiene límites precisos. Las imágenes no equivalen a las de un detector físico, cuya equivalencia no se demostró; su exposición es una escala convencional y no un tiempo calibrado de medición; los escenarios no reproducen un volcán concreto, y las pruebas realizadas no demuestran generalización a datos experimentales ni capacidad alguna de reconstrucción de densidad. Las realizaciones de baja estadística comparten referencias finitas que conservan su propio ruido, de modo que no deben presentarse como simulaciones independientes de alta estadística. La representatividad del conjunto se sostiene dentro de la familia idealizada, en los términos discutidos en el capítulo~\ref{chap:resultados}.

\section{Conclusión del segundo objetivo específico}
\label{sec:conc-oe2}

El OE2 se formuló en el plan de trabajo como «diseñar e implementar un modelo de aprendizaje profundo basado en redes generativas adversarias condicionales para transformar imágenes muográficas de baja estadística en imágenes de calidad comparable a las de alta estadística». Se diseñó e implementó una red generativa antagónica condicional de Wasserstein con penalización de gradiente, formada por un generador U-Net de 20\,543\,809 parámetros y un crítico PatchGAN condicionado por la entrada, para transformar imágenes muográficas con el 5\,\% de los conteos en estimaciones de la imagen de referencia de alta estadística. Antes de entrenarla, se corrigieron la coherencia entre la puntuación del crítico y su penalización de gradiente y la dependencia de la normalización respecto de la referencia, y se garantizó la reproducibilidad exacta de los entrenamientos interrumpidos. El modelo se entrenó con las doce geometrías de entrenamiento, se desarrolló mediante un piloto y una secuencia de ajustes decididos con validación, y se evaluó, junto con la U-Net entrenada con pérdida L1 y con seis continuaciones ajustadas, sobre tres geometrías de prueba y una de generalización, en una evaluación descriptiva sobre datos potencialmente expuestos.

En prueba, la cGAN aumentó el SSIM de 0,7643 a 0,8342 y el PSNR de 24,258 a 31,170~dB respecto de la entrada, y la U-Net alcanzó 0,8545 y 33,074~dB; en generalización se repitió el mismo patrón. Ambas redes, sin embargo, aumentaron el indicador regional de sesgo de cuentas de 1,832\,\% a 14,566 y 22,388\,\%, respectivamente. Los ajustes redujeron ese indicador a entre 3,996 y 5,825\,\% en prueba, a costa de una similitud menor que la de la U-Net, y el término espacial de borde redujo el error de borde entre un 31 y un 40\,\% en las cuatro geometrías reservadas, sin una mejora simultánea de todos los indicadores. Las doce comparaciones fijadas de antemano se clasificaron como compromiso o deterioro en al menos un indicador.

Con ello, los resultados sustentan el diseño, la implementación y el entrenamiento del modelo que el OE2 formuló, y su evaluación descriptiva con el 5\,\% de los conteos: el modelo reduce el ruido de conteo y aumenta la similitud con la referencia de manera consistente en escenarios no usados para entrenarlo. Según el SSIM y el PSNR, la cGAN alimentada con el 5\,\% de los conteos superó, en cada una de las cuatro geometrías reservadas, a la imagen sin procesar obtenida con el doble de conteos, aunque permaneció por debajo de la obtenida con diez veces más. Esa mejora de la similitud no garantiza la conservación de las cuentas regionales, que el uso muográfico de las imágenes exige considerar, de modo que la calidad comparable a la de alta estadística queda respaldada solo de forma parcial. La formulación del objetivo se refiere a la calidad de las imágenes y no establece la equivalencia física completa como condición independiente, que tampoco se añade aquí a posteriori ni se demostró. La conclusión se limita a referencias idealizadas, a un catálogo conocido cuyas escenas reservadas se habían evaluado antes y a la fracción de 0,05; no se extiende a volcanes reales, a la equivalencia física ni a otras fracciones de conteo.

\section{Conclusión del tercer objetivo específico}
\label{sec:conc-oe3}

El desempeño del modelo entrenado se evaluó con las métricas de similitud, de conservación de cuentas, de error de borde y de contraste de la evaluación final, y se comparó con la entrada sin red obtenida con el 1, el 5, el 10 y el 50\,\% de los conteos. Bajo la equivalencia idealizada entre fracción de conteos y exposición, la fracción de 0,05 corresponde a un factor nominal de reducción de 20 respecto de la referencia. Las redes alimentadas con el 5\,\% de los conteos superaron en SSIM y PSNR a la entrada sin red con el doble de conteos, cuya razón nominal de exposición es 2, y la cGAN lo hizo en cada una de las cuatro geometrías, y quedaron por debajo de la entrada con diez veces más conteos, salvo el PSNR de la U-Net en generalización, siempre con una peor conservación regional de las cuentas.

Esos resultados constituyen una comparación descriptiva y no la determinación de un factor de reducción alcanzable. No se estimó un factor entre las razones nominales de 2 y 10, un máximo ni un intervalo, porque las redes solo se evaluaron con una fracción, no hay mediciones intermedias y no se interpoló ni extrapoló entre los puntos de la rejilla. Determinar ese factor exige un criterio de calidad aceptable con un margen fundamentado, que no se había fijado; en su ausencia, el resultado se informa como una comparación descriptiva por métrica y, mientras no se establezca, el factor de reducción del tiempo de adquisición alcanzable con el esquema propuesto no está demostrado, y el valor nominal de 20 describe solo la reducción de conteos de la entrada.

\section{Conclusión general}
\label{sec:conc-general}

El proyecto se propuso desarrollar y validar un esquema de reconstrucción muográfica por transmisión asistido por aprendizaje profundo para reducir el tiempo efectivo de adquisición. Los tres objetivos específicos produjeron una cadena completa y verificable: un conjunto sintético de 9\,120 pares obtenido con Geant4 sobre diecinueve escenarios volcánicos idealizados, cuya integridad y reproducción desde los histogramas se comprobaron; un modelo generativo condicional cuya implementación, entrenamiento y evaluación se verificaron de forma independiente, y una comparación descriptiva de su desempeño frente a la exposición. Todos los productos quedaron identificados por sus huellas y acompañados de los protocolos y registros que permiten repetir su evaluación.

El resultado principal es la cuantificación de un compromiso. Con el 5\,\% de los conteos, el aprendizaje profundo mejora la similitud con la referencia más allá de lo que ofrece una adquisición sin procesar con el doble de conteos, pero lo hace desplazando las cuentas de las regiones físicamente relevantes, de modo que la mejora visual no equivale a una mejora de la información física de la imagen. Por ello, el esquema desarrollado no puede considerarse validado como un medio para reducir el tiempo de adquisición: esa validación requiere un criterio de calidad aceptable fijado de antemano, una evaluación confirmatoria sobre escenas no inspeccionadas y, para una aplicación real, la equivalencia con la respuesta de un detector físico, ninguno de los cuales estaba disponible. Dentro del alcance idealizado adoptado, el trabajo deja establecida la infraestructura experimental, la evidencia y los límites a partir de los cuales esa validación puede abordarse.

\section{Trabajo futuro}
\label{sec:trabajo-futuro}

Respecto del conjunto de datos, cuatro extensiones permitirían ampliar las afirmaciones del OE1 más allá del alcance idealizado. La primera es el control de equivalencia con un detector físico, que compara la imagen del estimador direccional con la de un panel expuesto a una fuente espacialmente amplia, con un criterio de equivalencia por intervalos simultáneos ya diseñado. La segunda es la incorporación de la topografía y de la estructura de densidades de un volcán concreto, en lugar de conos idealizados. La tercera es el transporte de partículas secundarias, del aire previo a la roca y del fondo instrumental, junto con la respuesta de un detector con parámetros fijados, que darían una escala de exposición física. La cuarta es una segunda ejecución independiente del transporte, que extendería la reproducción verificada desde el posproceso hasta la simulación.

Respecto del modelo y de su evaluación, el trabajo futuro se divide en mejoras metodológicas y en el camino hacia una aplicación de campo. Entre las primeras, la más inmediata es fijar, con una justificación externa o derivada de la aplicación, el criterio de calidad aceptable que separe la equivalencia descriptiva de la aceptación física, como condición para determinar un factor de reducción alcanzable. Con ese criterio, la evaluación debería repetirse sobre geometrías nuevas, generadas después de fijar el protocolo y nunca inspeccionadas, para obtener una evaluación confirmatoria, y extenderse a modelos entrenados con las demás fracciones de conteo, con repeticiones independientes desde cero que permitan estimar la variabilidad del entrenamiento. El compromiso observado sugiere, además, explorar formulaciones que restrinjan explícitamente la conservación de las cuentas, como pérdidas basadas en la verosimilitud del conteo o restricciones de la salida, y cuantificar la incertidumbre de las predicciones, junto con referencias de mayor exposición que acoten el ruido propio de la referencia. Hacia una aplicación de campo, el modelo debería entrenarse y evaluarse con imágenes que incorporen la respuesta de un detector físico, el fondo y una topografía real, y contrastarse finalmente con adquisiciones de exposición prolongada de un volcán concreto, que son las únicas que permitirían medir una reducción efectiva del tiempo de adquisición.
